% !TEX root = ACA2019_abstract_pdfgen.tex
% This file should be rename with the speaker's last name. (e.g : smith)

% Write the author names in the order you wish them to appear
% and underline the one who will give the talk.

% Only the speaker's email address is required.

\begin{abstract_online}{Dynamic Applications for Learning and Exploring Mathematics Using Computer Algebra}{%
        \underline{William Bauldr}y$^{1}$, \underline{Wade Ellis}$^{2}$}{[BauldryWC@appstate.edu]
        }{%
        $^1$ Mathematical Sciences, Appalachian State University, Boone, NC. \newline{}
        $^2$ Mathematics Department, West Valley College, Saratoga, California.%
        }
    
%%% Start abstract 
We discuss designing self-contained electronic documents that form a microworld for student investigations. The documents include a CAS application in which students engage in `sandboxed' mathematical exploration. The inquiry-based exploration is led by a set of questions in the document that guide students, experimenting in the computer algebra and dynamic geometry microworlds, that are formulated under the \emph{Action-Consequence-Reflection paradigm}.
%%% End abstract

\textbf{Keywords} \newline{} Dynamic computer algebra pedagogical applications, Action-Consequence-Reflection paradigm

\textbf{References} \newline{}
[1] \textsc{D.~Apple} and \textsc{W.~Ellis}, ``Learning how to learn: Improving the performance of learning,'' Int'l J of Process Education, 7(1), 2015, 21-28.

[2] \textsc{G.~Burrill}, ``The Role of Handheld Technology in Teaching and Learning Secondary School Mathematics,'' ICME 11, TSG-22. Monterrey, M\'{e}xico, 2008.

[3] \textsc{T.~Dick}, \textsc{G.~Burrill}, and \textsc{G.~Brady}, ``New technologies offer new ways to engage students,'' NCSM Newsletter, 38 (2007).

[4] \textsc{W.~Ellis}, ``Technology and Calculus.'' In \emph{Calculus Renewal: Issues for Undergraduate Mathematics Education in the Next Decade}, 53-68, S.~Ganter (ed), Springer US, Boston, 2000.

[5] \textsc{W.~Ellis, W.~Bauldry, M.~Boss\'e,} and \textsc{S, Oteh}, ``Employing Technology to Visualize Complex Roots of Real Polynomials,'' Electronic Proc.~of TIME-2016, UNAM, M\'{e}xico City, Jan.~2017.

[6] \textsc{W.~Bauldry}, ``Using Maple in Modern Algebra and Advanced Calculus Courses,'' Proc. of the 20th Annual ICTCM, 2009, 13-17.

\end{abstract_online}
